\begin{table}[!h]
\centering
\caption{\label{tab:data_sources}Data Sources}
\centering
\fontsize{8}{10}\selectfont
\begin{threeparttable}
\begin{tabular}[t]{lccc}
\toprule
Data Source & Coverage & Years & Key Variables\\
\midrule
\addlinespace[0.3em]
\multicolumn{4}{l}{\textit{\textbf{Panel A: Electoral Data}}}\\
\hspace{1em}Rajasthan SEC & 11,506 GPs & 2005, 2010, 2015, 2020 & Reservation status, winner sex\\
\hspace{1em}Uttar Pradesh SEC & 58,994 GPs & 2005, 2010, 2015 & Reservation status, winner sex\\
\hspace{1em}Candidate Affidavits & Rajasthan 2020 & 2020 & Age, education, assets\\
\addlinespace[0.3em]
\multicolumn{4}{l}{\textit{\textbf{Panel B: Geographic \& Socioeconomic Data}}}\\
\hspace{1em}Local Government Directory & India & 2020 & GP-village mapping\\
\hspace{1em}SHRUG 2.0 & India & 1991, 2001, 2011 & Village-level covariates\\
\hspace{1em}Census Village Directory & India & 2001 & Infrastructure, services\\
\addlinespace[0.3em]
\multicolumn{4}{l}{\textit{\textbf{Panel C: Primary Data}}}\\
\hspace{1em}Phone Audit (Quota) & Rajasthan & 2024 & 500 sampled, 377 answered\\
\hspace{1em}Phone Audit (Open) & Rajasthan & 2024 & 507 sampled, 400 answered\\
\bottomrule
\end{tabular}
\begin{tablenotes}
\item \textit{Note: } 
\item SEC = State Election Commission. GP = Gram Panchayat. SHRUG = Socioeconomic High-resolution Rural-Urban Geographic Platform for India. LGD = Local Government Directory.
\end{tablenotes}
\end{threeparttable}
\end{table}
